\documentclass[11pt,a4paper]{article}
\usepackage{notes}
\usepackage{fancyhdr}

\pagestyle{fancy}
\fancyhf{}
\lhead{Geophysics IIIA: Potential Fields \& Geothermics}
\rhead{Week 9 Plate Cooling}
\cfoot{\thepage}

\title{Week 9 --- Activity: Plate Cooling and Fourier Series}
\author{Geophysics 3A: Potential Fields \& Geothermics}
\date{\today}

\begin{document}
\maketitle

\vspace{1em}

\section*{Purpose}
To understand how imposing a fixed lithospheric thickness modifies transient
cooling and why Fourier series naturally arise in plate cooling models.

\section*{Concept}
Plate models introduce a basal boundary condition at a fixed depth, representing
the lithosphere--asthenosphere boundary.

\section*{Mathematical Form}
The temperature field can be expressed as a Fourier series:
\[
T(z,t) = \sum_{n=1}^{\infty} A_n
\sin\!\left(\frac{n\pi z}{L}\right)
\exp\!\left(-\frac{n^2 \pi^2 \kappa t}{L^2}\right).
\]

\section*{Tasks}
\begin{enumerate}
  \item Explain how the index $n$ relates to thermal wavelength.
  \item Why do short-wavelength modes decay faster than long-wavelength modes?
  \item Compare thermal evolution predicted by plate cooling and half-space cooling.
\end{enumerate}

\section*{Geological Interpretation}
\begin{itemize}
  \item Plate thickness limits lithospheric growth.
  \item Explains flattening of bathymetry and heat flow for old oceanic plates.
\end{itemize}

\section*{Connection}
Relate Fourier modes to spectral filtering concepts used previously in gravity
and magnetic field analysis.

\end{document}